\chapter{Model Architecture and Evaluation Methodology}
\label{ch:tools}

\section{Experimental Methodology and Objectives}

Establishing a rigorous modeling and evaluation methodology required resolving two primary technical questions:
1. \emph{Given constrained computational hardware (a single GPU instance), which model architecture and parameter-efficient fine-tuning method offer feasibility and training efficiency?}
2. \emph{What evaluation protocols are necessary to ensure that reported accuracy improvements represent genuine trajectory modeling rather than dataset artifacts?}

The selection of model architecture and fine-tuning mechanics remained fixed across experimental iterations, ensuring consistent cross-experiment comparisons. The second question led to three fundamental evaluation principles designed to prevent false positives and over-estimation of baseline performance gains.

\paragraph{Data Provenance.} Empirical measurements presented in this chapter derive from three early dataset cohorts documented in \textbf{Table~\ref{tab:datasets}}: raw logs (\Draw{}), basic filtered logs (\Dfilt{}), and active mobility logs (\Dmob{}). Each cohort evaluates models trained on \textbf{400 subscribers} and graded on \textbf{100 held-out subscribers}. Corresponding persistence baseline accuracies across these three datasets are 69.2\%, 79.0\%, and 45.5\%, respectively.

\section{Parameter-Efficient Fine-Tuning and Model Architecture}
\label{sec:tools-peft}

Fine-tuning pre-trained language models can be conducted under varying memory footprints. Fine-tuning the full parameter set requires substantial memory resources, whereas Parameter-Efficient Fine-Tuning (PEFT) methods freeze pre-trained parameters and train auxiliary low-rank matrices.

% F5b -- Full fine-tuning, LoRA and QLoRA, drawn side by side.
\begin{figure}[!htbp]
\centering
\fitwidth{%
\begin{tikzpicture}[font=\small,
  big/.style={draw=clrule,fill=clnonsig!30,minimum width=2.5cm,minimum height=2.5cm},
  small/.style={draw=clgain,fill=clgain!25,minimum width=2.5cm,minimum height=0.42cm},
  smallv/.style={draw=clgain,fill=clgain!25,minimum width=0.42cm,minimum height=2.5cm},
  q/.style={draw=clmodel,fill=clmodel!18,minimum width=2.5cm,minimum height=2.5cm},
  lbl/.style={font=\scriptsize,text=clrule,align=center},
]
% ---------------- (a) full fine-tuning ----------------
\node[big,fill=clloss!25,draw=clloss] (fa) at (0,0) {};
\node[lbl,above=3pt of fa] {\textbf{(a) Full fine-tuning}};
\node[lbl,below=3pt of fa,align=center] {every weight is updated\\ \textcolor{clloss}{100\,\% trainable}\\ does not fit on a free GPU};
\node[font=\scriptsize,text=clloss] at (fa) {$W$};

% ---------------- (b) LoRA ----------------
\node[big] (fb) at (4.6,0) {};
\node[font=\scriptsize,text=clrule] at (fb) {$W$ \emph{frozen}};
\node[smallv] (bA) at (6.35,0.55) {};
\node[small]  (bB) at (6.35,-1.0) {};
\node[font=\scriptsize] at (bA) {$A$};
\node[font=\scriptsize] at (bB) {$B$};
\node[lbl,above=3pt of fb] {\textbf{(b) LoRA}};
\node[lbl,below=17pt of fb,align=center,xshift=0.9cm]
  {the frozen matrix is left alone; only the two\\ thin \textcolor{clgain}{$A$} and
   \textcolor{clgain}{$B$} matrices are trained\\ \textcolor{clgain}{$\approx$2\,\% trainable}};

% ---------------- (c) QLoRA ----------------
\node[q] (fc) at (10.4,0) {};
\node[font=\scriptsize,text=clmodel,align=center] at (fc) {$W$ \emph{frozen}\\ \emph{and 4-bit}};
\node[smallv] (cA) at (12.15,0.55) {};
\node[small]  (cB) at (12.15,-1.0) {};
\node[font=\scriptsize] at (cA) {$A$};
\node[font=\scriptsize] at (cB) {$B$};
\node[lbl,above=3pt of fc] {\textbf{(c) QLoRA}, used here};
\node[lbl,below=17pt of fc,align=center,xshift=0.9cm]
  {same, plus the frozen matrix is stored\\ at 4 bits per number instead of 16\\
   \textcolor{clmodel}{$\approx$1/4 of the frozen memory}};
\end{tikzpicture}}

\caption{What LoRA and QLoRA do}
\label{fig:lora}
\figtext{%
\figpara{The grey square is the model.}{A language model is essentially a stack
of very large tables of numbers, and adapting one to a new task means changing
numbers in those tables. The three panels are the three ways of doing it, drawn
at the same scale.}

\figpara{(a) Change everything.}{Every number in the table is modified. This has
the highest ceiling in principle, but while the run lasts each number has to be
held in the graphics card's memory together with about three more used to keep
track of how to adjust it. The whole thing is several times larger than the
model, and it does not fit on a free card.}

\figpara{(b) LoRA: leave the table alone and add two thin ones beside
it.}{\cite{hu2022lora} The big table is frozen and never touched. Two narrow
tables are added next to it, and only those are trained. Two narrow tables
standing in for a change to a wide one is the whole of the idea, and here it
means about 2 numbers in every 100 are ever allowed to move.}

\figpara{(c) QLoRA: the same, with the frozen table compressed.}{\cite{dettmers2023qlora}
This is the method used throughout the report. On top of LoRA, the frozen table
is stored with four bits per number instead of sixteen, which quarters the memory
it occupies. The compression is safe precisely because those numbers are frozen:
they are read, never updated. The added tables stay at full precision.}

\figpara{Why this decides the shape of the whole internship.}{Not accuracy, but
waiting time. With (c), one training run fits in half an hour on a free graphics
card, so nine weeks buy thirty experiments instead of two.}}

\end{figure}


\textbf{Figure~\ref{fig:lora}} compares parameter adaptation mechanisms across three memory scales: full fine-tuning (\textbf{\panelref{fig:lora}{a}}), Low-Rank Adaptation, LoRA (\textbf{\panelref{fig:lora}{b}}), and Quantized Low-Rank Adaptation, QLoRA (\textbf{\panelref{fig:lora}{c}}). Full parameter updating (\textbf{\panelref{fig:lora}{a}}) modifies 100\% of model parameters, requiring memory to maintain optimizer states for all weights and exceeding the VRAM capacity of standard single-GPU instances. LoRA (\textbf{\panelref{fig:lora}{b}}) \textbf{\cite{hu2022lora}} freezes the original base weight matrices $W$ and inserts auxiliary low-rank decomposition matrices $A$ and $B$, training only $\approx 2\%$ of total parameters. QLoRA (\textbf{\panelref{fig:lora}{c}}) \textbf{\cite{dettmers2023qlora}} further compresses frozen base parameters into 4-bit NormalFloat precision (reducing base memory footprint to $\approx 1/4$) while training 16-bit LoRA adapter weights. QLoRA was selected to enable model execution within 16\,GB GPU memory limits. This selection primarily impacts training throughput rather than model accuracy, accelerating execution from 5.5 hours to 30 minutes per training run and supporting extensive hyperparameter sweeps and per-group adapter routing experiments (Chapter~\ref{ch:groups}).

\subsection{Model Selection and Capacity Evaluation}
\label{sec:tools-model}

Initial experiments evaluated \textbf{Mistral-7B} \textbf{\cite{jiang2023mistral}} (7 billion parameters). A preliminary 30-minute test run on 1.2\% of the dataset yielded 16.87\% top-1 accuracy, verifying pipeline end-to-end functionality. (Comparative tests using Llama models failed to produce syntactically valid cell identifier output.)

However, full training of Mistral-7B required 5.5 hours per run on a single T4 GPU, reaching a maximum accuracy of 41.1\%. In contrast, \textbf{Qwen2.5-0.5B-Instruct} \textbf{\cite{qwen2025}} (500 million parameters, 14$\times$ smaller) completed 49 training epochs in 72 minutes (reducing to 29 minutes under optimized early stopping). Consequently, Qwen2.5-0.5B was chosen for all subsequent experiments.

Because next-cell prediction across 369 discrete vocabulary items requires modeling localized spatial transition patterns rather than general world knowledge or natural language reasoning, the smaller model offers adequate capacity. The $14\times$ parameter reduction provided a $23\times$ speedup per run, enabling extensive experimental iteration within computational constraints.

% F5a -- The tools, as a stack: what sits on top of what, and who does what.
\begin{figure}[!htbp]
\centering
\fitwidth{%
\begin{tikzpicture}[font=\small,
  lay/.style={draw=clrule,rounded corners=2pt,align=center,inner sep=6pt,
              minimum height=1.35cm,text width=11.4cm,anchor=center},
  box/.style={draw=clrule,rounded corners=2pt,align=center,inner sep=6pt,
              minimum height=1.75cm,text width=3.35cm,font=\footnotesize,anchor=center},
  wide/.style={box,text width=5.35cm},
  side/.style={draw=clrule,fill=white,rounded corners=2pt,align=center,
               inner sep=6pt,text width=3.3cm,font=\footnotesize,anchor=center},
  tag/.style={font=\scriptsize,text=clrule,anchor=east},
]
% --- layer 4: my own code -------------------------------------------------
\node[lay,fill=clmodel!12,draw=clmodel] (mine) at (0,4.75)
  {\textbf{My notebook}\\[2pt]\footnotesize
   the custom training loop: which mistakes count more, which set of numbers
   answers for which kind of user, and the whole grading};
\node[tag] at (-6.1,4.75) {written by me};

% --- layer 3: the three libraries ----------------------------------------
\node[box,fill=clsoft] (tr)   at (-3.95,2.55)
  {\textbf{Transformers}\\[2pt] loads the model and runs a standard training
   loop};
\node[box,fill=clsoft] (peft) at (0,2.55)
  {\textbf{PEFT}\\[2pt] freezes the model and adds the small trainable part};
\node[box,fill=clsoft] (trl)  at (3.95,2.55)
  {\textbf{TRL}\\[2pt] the format the training file is written in};
\node[tag] at (-6.1,2.55) {the libraries};

% --- layer 2: model + 4-bit storage --------------------------------------
\node[wide,fill=clink!14] (qwen) at (-2.95,0.35)
  {\textbf{Qwen2.5-0.5B}\\[2pt] the model itself, 500 million numbers,
   downloaded already trained on general text};
\node[wide,fill=clink!14] (bnb) at (2.95,0.35)
  {\textbf{bitsandbytes}\\[2pt] stores those frozen numbers in 4 bits each
   instead of 16, so they fit};
\node[tag] at (-6.1,0.35) {what is trained};

% --- layer 1: the machine -------------------------------------------------
\node[lay,fill=clnonsig!35] (colab) at (0,-1.65)
  {\textbf{Google Colab, free tier: one NVIDIA T4 graphics card, 16\,GB}\\[2pt]
   \footnotesize rented in the browser, a few hours a day, and liable to
   disconnect};
\node[tag] at (-6.1,-1.65) {where it runs};

% --- the side column: analysis -------------------------------------------
\node[side,anchor=west] (an) at (6.55,1.45)
  {\textbf{pandas}\\ \textbf{scikit-learn}\\ \textbf{scipy}\\[4pt]
   preparing the files, splitting people into train and test, and checking the
   split is fair.\\[3pt] No modelling here.};

\draw[clrule,dashed] (6.15,-2.45) -- (6.15,5.55);
\end{tikzpicture}}
\caption{The tools, and the job of each one}
\label{fig:toolchain}

\end{figure}


\textbf{Figure~\ref{fig:toolchain}} illustrates the four-layer software and hardware stack utilized throughout this study:
- \textbf{Hardware Execution Layer} (bottom): Google Colab single NVIDIA T4 GPU instances (16\,GB VRAM), establishing the primary computational memory constraint.
- \textbf{Model & Quantization Layer}: Qwen2.5-0.5B base parameters quantized to 4-bit precision using \texttt{bitsandbytes}.
- \textbf{Open-Source Library Layer}: Hugging Face \texttt{Transformers} for model loading and optimization loops; \texttt{PEFT} for managing LoRA adapter parameters; and \texttt{TRL} for sequence dataset formatting.
- \textbf{Custom Modeling Layer} (top): custom Python training and evaluation routines implementing weighted cross-entropy loss, dynamic sequence sampling, per-group LoRA adapter routing, and evaluation protocols.
- \textbf{Analytical Tooling} (right): standard scientific computing libraries (\texttt{pandas}, \texttt{scikit-learn}, \texttt{scipy}) managing offline data splitting, feature engineering, and statistical bootstrap resampling without touching neural network parameters.

Reaching inside the training step to implement custom loss weighting and parameter routing represented the primary design requirement governing the choice of this software toolchain.

\section{Data Representation and Tokenization Scheme}

Language models process discrete token sequences. \textbf{Figure~\ref{fig:training-example}} illustrates the formatting of a subscriber daily trajectory as it is stored on disk (\textbf{\panelref{fig:training-example}{a}}) and as it is handed to the language model tokenizer (\textbf{\panelref{fig:training-example}{b}}).

% F3b -- One training example: the JSONL record, and the same line as tokens.
\begin{figure}[!htbp]
\centering

{\textbf{\panelnum{fig:training-example}{a} One person's day, as it is stored on disk.} One line of the
training file, spread over several lines here so that it can be read. Nothing is
a conversation in any real sense: the three labels are simply the slots the
training library expects to find.\par}
\vspace{3pt}

\begin{lstlisting}
{"conversations": [
  {"role": "system",    "content": "You are an AI that generates a user's
                                    cell-visit trajectory."},
  {"role": "user",      "content": "Here is the sequence for this user."},
  {"role": "assistant", "content": "User 2832452 | 200 events |
                                    BSOCHO1 BSOCHO1 UKVCHE102 BSOCHO1 ..."}
]}
\end{lstlisting}

\vspace{2pt}
{\footnotesize
\begin{tabular}{@{}p{2.2cm}p{10.75cm}@{}}
\texttt{system} & What the task is. The same words on every line of the file. \\
\texttt{user}   & The request. Also the same everywhere: a placeholder, not a
                  question. \\
\texttt{assistant} & The person's actual day. This is what the model has to
                  learn to produce, and the only part its mistakes are counted
                  on. \\
\end{tabular}\par}

\vspace{0.55cm}
\hrule height 0.4pt
\vspace{0.45cm}

{\textbf{\panelnum{fig:training-example}{b} The same line, once it has been cut into symbols.} This is what
the model actually reads. Each antenna name is \emph{one} symbol, not a spelling;
each event is preceded by its hour. Only the dark blue positions are graded.\par}
\vspace{5pt}

\fitwidth{%
\begin{tikzpicture}[font=\scriptsize,node distance=1.2pt,
  tok/.style={draw=clrule,minimum height=0.55cm,inner xsep=3.5pt,inner ysep=0pt},
  hr/.style={tok,draw=clmove,fill=clmove!20},
  cl/.style={tok,draw=clmodel,fill=clmodel!18},
  pfx/.style={tok,fill=clsoft,text=clrule},
]
\node[pfx] (t1) {\textless system\textgreater};
\node[pfx,right=of t1] (t2) {\dots};
\node[pfx,right=of t2] (t3) {User 2832452};
\node[pfx,right=of t3] (t4) {\textbar};
\node[hr,right=of t4]  (t5) {\textless H07\textgreater};
\node[cl,right=of t5]  (t6) {BSOCHO1};
\node[hr,right=of t6]  (t7) {\textless H09\textgreater};
\node[cl,right=of t7]  (t8) {BSOCHO1};
\node[hr,right=of t8]  (t9) {\textless H12\textgreater};
\node[cl,right=of t9]  (t10) {UKVCHE102};
\node[hr,right=of t10] (t11) {\textless H18\textgreater};
\node[cl,right=of t11] (t12) {BSOCHO1};
\node[pfx,right=of t12] (t13) {\dots};

\draw[decorate,decoration={brace,amplitude=4pt,mirror},clrule]
  ($(t1.south west)+(0,-2pt)$) -- ($(t4.south east)+(0,-2pt)$)
  node[midway,below=5pt,font=\scriptsize,text=clrule,align=center] {opening\\ \emph{not graded}};
\draw[decorate,decoration={brace,amplitude=4pt,mirror},clmodel]
  ($(t5.south west)+(0,-2pt)$) -- ($(t13.south east)+(0,-2pt)$)
  node[midway,below=5pt,font=\scriptsize,text=clmodel,align=center]
  {the day itself: only the 4 blue positions are graded};
\end{tikzpicture}}

\vspace{0.15cm}
\vspace{-0.15cm}
\replegend{\swatch{clsoft}~opening: excluded from the grade \qquad
           \swatch{clmove!20}~hour: \emph{given}, never guessed \qquad
           \swatch{clmodel!18}~antenna name: the only positions graded}

\caption{One person's day, as the model receives it}
\label{fig:training-example}

\end{figure}


As shown in \textbf{\panelref{fig:training-example}{a}}, daily trajectories are structured on disk as JSON Lines records containing system, user, and assistant role slots to conform with \texttt{TRL} library conventions. When tokenized in \textbf{\panelref{fig:training-example}{b}}, three specific encoding choices are implemented:
1. \textbf{Dedicated Vocabulary Tokens for Cell IDs}: All 369 cell identifiers were added directly to the model's vocabulary as atomic tokens (dark blue). This prevents subword tokenization splitting strings like \cid{BSOCHO1} into sub-tokens, forcing atomic cell predictions.
2. \textbf{Explicit Hourly Tokens}: Temporal context is represented by inserting explicit hour tokens (\cid{<H00>} through \cid{<H23>}, orange) preceding each cell event. Hour tokens are provided as inputs for free during inference and consume one token position per event (compared to 5 tokens for text strings like ``07:31'').
3. \textbf{Loss Masking}: Cross-entropy loss is computed exclusively on cell identifier token positions (dark blue), masking header boilerplate (grey) and temporal tokens.

\subsection{Training Pipeline and Optimization Mechanics}

% F5d -- How one training run actually works, step by step.
\begin{figure}[!htbp]
\centering
\fitwidth{%
\begin{tikzpicture}[font=\small,
  st/.style={draw=clrule,fill=clsoft,rounded corners=3pt,align=center,
             inner sep=6pt,text width=2.85cm,minimum height=2.75cm},
  ar/.style={-{Stealth[length=5pt]},clrule,very thick},
  num/.style={circle,draw=clmodel,fill=white,inner sep=1.2pt,
              font=\scriptsize\bfseries,text=clmodel},
]
\node[st] (a) at (0,0)
  {{\footnotesize\bfseries\color{clmodel} Take a slice}\\[2pt]\footnotesize
   a stretch of one training person's day: hour, antenna, hour, antenna};
\node[st] (b) at (3.55,0)
  {{\footnotesize\bfseries\color{clmodel} Hide the next antenna}\\[2pt]\footnotesize
   the model reads the slice and names what it thinks comes next};
\node[st] (c) at (7.10,0)
  {{\footnotesize\bfseries\color{clmodel} Measure the surprise}\\[2pt]\footnotesize
   compare with the antenna really visited; a confident wrong answer costs more
   than a hesitant one};
\node[st] (d) at (10.65,0)
  {{\footnotesize\bfseries\color{clmodel} Nudge the small part}\\[2pt]\footnotesize
   only the added numbers move; the 500 million original ones never do};

\draw[ar] (a) -- (b); \draw[ar] (b) -- (c); \draw[ar] (c) -- (d);
\node[num] at (a.north west) {1}; \node[num] at (b.north west) {2};
\node[num] at (c.north west) {3}; \node[num] at (d.north west) {4};

% --- the loop back --------------------------------------------------------
\draw[ar] (d.south) -- ++(0,-0.65) -- ++(-10.65,0) -- (a.south);
\node[font=\scriptsize,text=clrule] at (5.32,-2.42)
  {repeat, slice after slice, until every training person has been read once:
   that is \textbf{one pass}};

% --- the checkpoint rule --------------------------------------------------
\node[draw=clgain,fill=clgain!10,rounded corners=3pt,align=center,inner sep=6pt,
      text width=13.4cm,font=\small,anchor=north] (ck) at (5.32,-2.85)
  {{\footnotesize\bfseries\color{clgain} 5. At the end of every pass: grade the
    model on people it is not training on, save a copy, and keep only the best
    copy of the whole run}\\[2pt]\footnotesize
   On this task the best copy arrives between the fourth and the ninth pass, and
   everything after it is worse. Keeping the last copy instead of the best one
   costs 15.5 correct answers per hundred (Chapter~\ref{ch:first}).};
\end{tikzpicture}}
\caption{How one training run works}
\label{fig:training-loop}
\figtext{%
\figpara{Steps 1 to 4 are the whole of training.}{A slice of somebody's day is
shown to the model with the next antenna hidden. The model names an antenna. The
answer is compared with the truth and turned into a single number measuring how
surprised the model was, so that being confidently wrong costs more than
hesitating. That number is then used to nudge the model's adjustable numbers
very slightly in the direction that would have made it less wrong. Then the next
slice.}

\figpara{Only a small part of the model ever moves.}{Step 4 is where the method
of Figure~\ref{fig:lora} bites. The 500 million numbers the model was published
with are frozen and never touched; the nudging applies only to a small set of
numbers added beside them, about 2 out of every 100. That is what makes a
training run fit in half an hour on a free graphics card.}

\figpara{Step 5 is not an optimisation, it is a safety rule.}{Left to itself, a
model trained on a few hundred short days stops learning the habits of the
population and starts learning the days themselves by heart. When that happens
its score on people it has never seen turns around and falls, while its score on
its own training data keeps improving, so the training figures alone will happily
report success for another forty passes. Grading on separate people after each
pass, and keeping the best copy rather than the last, is what prevents that, and
it was the largest single improvement of the whole internship.}}
\end{figure}


\textbf{Figure~\ref{fig:training-loop}} details the five-stage training optimization loop:
1. \textbf{Sequence Slicing}: sampling sub-sequence windows of timestamped cell events.
2. \textbf{Target Masking & Forward Pass}: predicting the next cell token given preceding historical context.
3. \textbf{Surprise Loss Computation}: calculating cross-entropy loss between model prediction probabilities and true target cells.
4. \textbf{Parameter Gradient Updates}: computing backpropagation gradients strictly for trainable LoRA adapter weights (~2\% of parameters) while base parameters remain frozen.
5. \textbf{Validation Checkpointing & Early Stopping}: evaluating model accuracy on unseen held-out validation subscribers after each training pass and saving the optimal checkpoint.

Two additions were introduced from the sixth experimental iteration onward:
- \textbf{Weighted Cross-Entropy Loss}: Loss terms at specific hourly positions are scaled by a multiplier (e.g., $4\times$ multiplier for 18h--20h transitions) at Step 3 of \textbf{Figure~\ref{fig:training-loop}}, emphasizing critical transition errors without duplicating data.
- \textbf{Targeted Sequence Sampling}: Training slices containing key transition windows are sampled up to $3\times$ more frequently during batch generation at Step 1 of \textbf{Figure~\ref{fig:training-loop}}.

\subsection{Hardware Environment and Computational Resources}

All models were trained on \textbf{Google Colab} GPU instances equipped with single NVIDIA \textbf{T4 GPUs} (16\,GB VRAM). Standard scientific computing libraries (\texttt{pandas}, \texttt{scikit-learn}, \texttt{scipy}) managed dataset preprocessing, stratified splits, and statistical bootstrap resampling \textbf{\cite{efron1993bootstrap}}.

\subsection{Integration of Peer Contributions}

Data extraction tooling for generating daily trajectory logs and active mobility cohorts was developed by Nino. The framing of active mobility trajectory filtering was developed by Tiphain and Benjamin. Trajectory length equalisation techniques evaluated in Chapter~\ref{ch:scale} adapt methods formulated by Tiphain.

\section{Evaluation Protocol and Metrics}
\label{sec:scoring}

Model evaluation follows a standardized protocol: each test trajectory is partitioned into a historical prefix (defaulting to 80\% of the daily trajectory) handed to the model, and a target suffix (20\%) evaluated sequentially. At each target step, the model predicts the next cell given actual past history, preventing autoregressive error accumulation.

% F5 -- Where a day is cut, and how a single position is scored.
\begin{figure}[!htbp]
\centering
\fitwidth{%
\begin{tikzpicture}[font=\small,
  ar/.style={-{Stealth[length=4pt]},clrule},
  ctx/.style={draw=clrule,fill=clsoft,rounded corners=2pt,inner sep=4pt,
              font=\scriptsize,text width=3.1cm,align=center},
  prd/.style={draw=clmodel,fill=clmodel!8,rounded corners=2pt,inner sep=4pt,
              font=\scriptsize,text width=2.2cm,align=center},
]
% ---- (a) the split ------------------------------------------------------
\node[anchor=west,font=\footnotesize\bfseries,text=clmodel] at (0,1.05)
  {(a) Each person's day is cut in two, once and for all};
\fill[clmodel!22] (0,0) rectangle (11.1,0.65);
\fill[clink!25]  (11.1,0) rectangle (13.9,0.65);
\draw[clrule] (0,0) rectangle (13.9,0.65);
\node[font=\footnotesize] at (5.55,0.33)
  {\textbf{THE PAST}, 80\,\% of the day, given to the model};
\node[font=\footnotesize] at (12.5,0.33) {\textbf{THE REST}, 20\,\%};
\node[font=\scriptsize,anchor=north,text=clrule] at (6.95,-0.10)
  {where this cut falls is called the \emph{context fraction}; several
   experiments move it from 30\,\% to 90\,\%};

% ---- (b) scoring one position at a time ---------------------------------
\node[anchor=west,font=\footnotesize\bfseries,text=clmodel] at (0,-1.20)
  {(b) The rest is then answered one moment at a time, never in one go};
\begin{scope}[yshift=-1.60cm]
  \node[ctx] (b0) at (2.6,-0.35)  {given events 1 to 40};
  \node[prd] (p0) at (5.9,-0.35)  {name event 41};
  \node[ctx] (b1) at (2.6,-1.10)  {given events 1 to 41};
  \node[prd] (p1) at (5.9,-1.10)  {name event 42};
  \node[ctx] (b2) at (2.6,-1.85)  {given events 1 to 42};
  \node[prd] (p2) at (5.9,-1.85)  {name event 43};
  \draw[ar] (b0) -- (p0); \draw[ar] (b1) -- (p1); \draw[ar] (b2) -- (p2);
  \node[font=\scriptsize,anchor=west,text width=6.6cm,align=left] at (7.6,-1.10)
    {Each time, the model is handed the person's \emph{real} past, not its own
     previous answers. So it can neither see the future nor go astray on its own
     earlier mistakes. The one-line rule answers exactly the same moments, and
     simply repeats the last antenna in the history it was given.};
\end{scope}
\end{tikzpicture}}
\caption{Where the day is cut, and how each moment is answered}
\label{fig:protocol}
\figtext{%
\figpara{Panel (a) is the cut; panel (b) is how each moment is answered.}{Each
test person's day is split once into a part the model may read and a part it has
to produce, the default being 80/20. Where that cut falls is part of the
\emph{measurement} rather than a property of the model, which is why it is drawn
here (Chapter~\ref{ch:first}). The model is then not asked to invent a whole
afternoon: it is asked for one antenna, told what really happened, and asked for
the next, so it can neither see the future nor compound its own errors. The
one-line rule is run through exactly the same moments with exactly the same
history and simply answers the last antenna it was shown, which is what makes the
two numbers comparable at all.}}
\end{figure}


\textbf{Figure~\ref{fig:protocol}} illustrates this evaluation protocol: \textbf{\panelref{fig:protocol}{a}} shows the trajectory partition into an 80\% prefix context (given to the model) and a 20\% suffix target (evaluated). As depicted in \textbf{\panelref{fig:protocol}{b}}, suffix evaluation operates sequentially: at each target step, the model is queried for the single next cell, graded against ground truth, and then provided with the true historical event for the next prediction step. This prevents compounding autoregressive errors and ensures that the model and the persistence baseline are evaluated on identical historical sequences.

% F5e -- What top-1, top-3 and top-5 mean, and how a score is computed from them.
\begin{figure}[!htbp]
\centering
\fitwidth{%
\begin{tikzpicture}[font=\small,
  cand/.style={draw=clrule,fill=white,rounded corners=2pt,align=center,
               inner sep=4pt,text width=2.15cm,font=\scriptsize},
  hit/.style={cand,draw=clgain,fill=clgain!14,thick},
  brk/.style={clmodel,thick,|-|},
  cell/.style={draw=clrule,minimum width=0.85cm,minimum height=0.55cm,
               inner sep=0pt,font=\scriptsize},
  ok/.style={cell,draw=clgain,fill=clgain!40},
  ko/.style={cell,draw=clnonsig,fill=white},
]
% ================= (a) one prediction, ranked =================
\node[anchor=west,font=\footnotesize\bfseries,text=clmodel] at (-0.35,3.55)
  {(a) The model never answers once. It ranks all 369 antennas.};
\node[anchor=west,font=\scriptsize,text=clrule] at (-0.35,3.18)
  {Here are its first five, for one moment of one person's day, with how sure it
   is of each.};

% the three depth markers, stacked above the row
\draw[brk] (-0.05,2.62) -- (2.35,2.62);
\node[font=\scriptsize,text=clmodel,anchor=west] at (2.55,2.62) {\textbf{top-1}};
\draw[brk] (-0.05,2.20) -- (7.45,2.20);
\node[font=\scriptsize,text=clmodel,anchor=west] at (7.65,2.20) {\textbf{top-3}};
\draw[brk] (-0.05,1.78) -- (12.55,1.78);
\node[font=\scriptsize,text=clmodel,anchor=west] at (12.75,1.78) {\textbf{top-5}};

\node[cand] (k1) at (1.15,0.95) {\#1\\ \cid{UKVBER101}\\ 52\,\%};
\node[cand] (k2) at (3.70,0.95) {\#2\\ \cid{UKVTGM104}\\ 23\,\%};
\node[hit]  (k3) at (6.25,0.95) {\#3\\ \cid{BKVPAN1}\\ 13\,\%};
\node[cand] (k4) at (8.80,0.95) {\#4\\ \cid{UKVDVO104}\\ 7\,\%};
\node[cand] (k5) at (11.35,0.95){\#5\\ \cid{BKVBOH1}\\ 5\,\%};
\node[font=\scriptsize,text=clrule,anchor=west] at (12.75,0.95)
  {\dots\ 364 more};

\node[font=\scriptsize,text=clgain,anchor=north,align=center] at (6.25,0.28)
  {$\uparrow$ this is the antenna the person really visited};

\node[font=\scriptsize,anchor=west] at (-0.35,-0.55)
  {\textcolor{clloss}{\textbf{top-1: wrong}}, it is not first. \qquad
   \textcolor{clgain}{\textbf{top-3: right}}, it is in the first three. \qquad
   \textcolor{clgain}{\textbf{top-5: right}}, it is in the first five.};

% ================= (b) how a score is computed =================
\draw[clrule,dotted] (-0.45,-1.05) -- (14.4,-1.05);
\node[anchor=west,font=\footnotesize\bfseries,text=clmodel] at (-0.35,-1.45)
  {(b) The score is then just a count, over every graded moment of every test
   person.};

\foreach \i/\st in {0/ok,1/ko,2/ok,3/ok,4/ko,5/ok,6/ok,7/ok,8/ko,9/ok}{
  \node[\st] at (0.1+\i*0.92,-2.25) {};
}
\node[font=\scriptsize,anchor=west,text width=6.2cm] at (9.4,-2.25)
  {shaded = right, empty = wrong\\[2pt]
   7 right out of 10 graded moments\\ $\Rightarrow$ a score of \textbf{70\,\%}};
\node[font=\scriptsize,text=clrule,anchor=west] at (-0.35,-2.95)
  {The one-line rule is graded on the very same moments, and its own count is
   what my score is compared against.};
\end{tikzpicture}}
\caption{Top-1, top-3 and top-5, and the way a score is obtained}
\label{fig:topk}
\figtext{%
\figpara{Panel (a): the model does not give an answer, it gives an order.}{Asked
what comes next, it assigns a degree of confidence to all 369 antennas and sorts
them, most likely first: here 52\,\% for \cid{UKVBER101}, 23\,\% for
\cid{UKVTGM104}, and so on. \textbf{Top-1} asks whether the right answer is in
first place, \textbf{top-3} whether it is anywhere in the first three,
\textbf{top-5} in the first five. In the example the right answer sits third, so
the same prediction fails at top-1 and succeeds at the other two. Reporting three
depths is the convention of this literature rather than a choice of mine
\cite{luca2021survey,qin2025mobllm}.}

\figpara{Panel (b): where a percentage comes from.}{The score is a plain count.
Every graded moment is answered, marked right or wrong at the chosen depth, and
the score is the number of right answers over the number of moments: ten moments,
seven right, 70\,\%. The one-line rule answers the very same moments and gets its
own count, and every result in this report is the difference between the two.}}

\end{figure}


\textbf{Figure~\ref{fig:topk}} demonstrates the multi-depth evaluation procedure. As shown in \textbf{\panelref{fig:topk}{a}}, the model outputs a probability distribution over all 369 candidate cell tokens in the vocabulary. \textbf{Top-1} accuracy evaluates whether the true target cell receives the highest probability score; \textbf{Top-3} and \textbf{Top-5} evaluate whether the true cell is present within the top three or top five ranked candidates, respectively. Reporting multi-depth accuracy follows standard mobility modeling literature conventions \textbf{\cite{luca2021survey,qin2025mobllm}}. As illustrated in \textbf{\panelref{fig:topk}{b}}, accuracy metrics represent simple percentage counts of correctly matched target steps divided by total evaluated steps across held-out test subscribers.

\section{Evaluation Methodology Principles}

% F -- The three rules about measuring, one panel each. A schematic: the only
% numbers printed on it are the three quoted in the text.
\begin{figure}[!htbp]
\centering
\fitwidth{%
\begin{tikzpicture}[font=\small,
  ttl/.style={anchor=north west,align=left,text width=4.5cm,
              font=\footnotesize\bfseries,text=clmodel},
  take/.style={anchor=north west,align=left,text width=4.34cm,inner sep=5pt,
               font=\scriptsize,fill=clsoft,draw=clrule,rounded corners=2pt},
  axl/.style={font=\tiny,text=clrule},
]
% ======================= panel (a) : read the peak =======================
\begin{scope}[shift={(0,0)}]
\node[ttl] at (0,0)
  {(a) Never read the last pass\\ of a run};
\draw[clrule] (0.25,-1.05) -- (0.25,-3.55) -- (4.45,-3.55);
\node[axl,rotate=90,anchor=south] at (0.05,-2.30) {score on new people};
\node[axl,anchor=north] at (2.35,-3.70) {passes over the training data};
\draw[clmodel,very thick]
  plot[smooth,tension=0.7] coordinates
  {(0.45,-3.25) (0.90,-1.90) (1.30,-1.35) (1.55,-1.25)
   (2.10,-1.60) (2.85,-1.95) (3.60,-2.20) (4.30,-2.35)};
\draw[clrule,dashed] (0.25,-1.25) -- (4.30,-1.25);
\draw[clrule,dashed] (0.25,-2.35) -- (4.30,-2.35);
\node[star,star points=5,star point ratio=2.2,fill=clgain,draw=none,
      inner sep=1.6pt] at (1.55,-1.25) {};
\node[text=clloss,font=\footnotesize\bfseries] at (4.30,-2.35) {$\times$};
\draw[{Stealth[length=4pt]}-{Stealth[length=4pt]},clgain,thick]
  (3.15,-2.35) -- (3.15,-1.25);
\node[anchor=west,font=\scriptsize\bfseries,text=clgain] at (3.22,-1.80) {15.5};
\node[take] at (0.08,-4.28)
  {The score on unseen people peaks early, then decays for tens of passes.
   Reading the \textcolor{clloss}{end} of the run instead of its
   \textcolor{clgain}{best point} throws away 15.5 correct answers per
   hundred.};
\end{scope}

% ======================= panel (b) : the cut is the ruler ================
\begin{scope}[shift={(5.10,0)}]
\node[ttl] at (0,0)
  {(b) Where the day is cut\\ belongs to the ruler};
\draw[clrule] (0.25,-1.05) -- (0.25,-3.55) -- (4.45,-3.55);
\node[axl,rotate=90,anchor=south] at (0.05,-2.30) {score};
\node[axl,anchor=north] at (2.35,-3.70)
  {share of the day given: 30\,\% $\rightarrow$ 90\,\%};
\fill[clgain!25] (0.55,-3.05) -- (4.25,-1.50) -- (4.25,-1.20) -- (0.55,-2.75) -- cycle;
\draw[clbase,very thick]  (0.55,-2.75) -- (4.25,-1.20);
\draw[clmodel,very thick] (0.55,-3.05) -- (4.25,-1.50);
\node[anchor=south,rotate=22.7,font=\scriptsize,text=clbase]
      at (2.40,-1.93) {the one-line rule};
\node[anchor=north,rotate=22.7,font=\scriptsize,text=clmodel]
      at (2.40,-2.32) {my model};
\node[take] at (0.08,-4.28)
  {Giving the model more history raises its score, and raises the rule by the
   same amount, because it also moves what counts as ``the last antenna
   seen''. A rising score across cuts says nothing about the model.};
\end{scope}

% ======================= panel (c) : leave room ==========================
\begin{scope}[shift={(10.20,0)}]
\node[ttl] at (0,0)
  {(c) A test the one-line rule\\ already passes};
\draw[clrule] (0.35,-3.55) -- (4.45,-3.55);
\node[axl,rotate=90,anchor=south] at (0.15,-2.30) {correct answers per 100};
% bar A : a raw day
\fill[clsoft,draw=clrule] (0.85,-3.55) rectangle (1.95,-1.05);
\fill[clbase] (0.85,-3.55) rectangle (1.95,-1.82);
\node[font=\scriptsize\bfseries,text=white] at (1.40,-2.68) {69.2};
\node[font=\scriptsize,text=clgain] at (1.40,-1.44) {30.8};
\node[anchor=north,align=center,font=\scriptsize,text=clrule] at (1.40,-3.62)
  {raw day file};
% bar B : people who move
\fill[clsoft,draw=clrule] (2.85,-3.55) rectangle (3.95,-1.05);
\fill[clbase] (2.85,-3.55) rectangle (3.95,-2.41);
\node[font=\scriptsize\bfseries,text=white] at (3.40,-2.98) {45.5};
\node[font=\scriptsize,text=clgain] at (3.40,-1.73) {54.5};
\node[anchor=north,align=center,font=\scriptsize,text=clrule] at (3.40,-3.62)
  {people who move};
\node[take] at (0.08,-4.28)
  {Blue is what one line of code already takes; the pale part above it is what
   a model could still win. Choosing whom to measure on nearly doubled that
   room, from 30.8 to 54.5.};
\end{scope}
\end{tikzpicture}}
\caption{The three rules about measuring, and the worth of each}
\label{fig:three-rules}
\end{figure}


\subsection{Principle 1: Early Stopping and Overfitting Mitigation}

% F6 -- Four runs, four datasets, the same shape: a plateau under the baseline.
\begin{figure}[t]
\centering
\plotlink{\nbllm}{\fitwidth{%
\begin{tikzpicture}
\begin{groupplot}[
  group style={group size=2 by 2, horizontal sep=1.6cm, vertical sep=2.0cm,
               xlabels at=edge bottom, ylabels at=edge left},
  rep, width=0.46\linewidth, height=4.3cm,
  xlabel={epoch}, ylabel={validation top-1 (\%)},
  ymin=0, ymax=100, ytick={0,20,...,100},
  no markers, title style={yshift=2pt},
]
% ---------- (a) D1, unfiltered day, 500 users, 49 epochs ----------
\nextgroupplot[title={\footnotesize (a) \sraw, 49 epochs},
               xmin=0, xmax=50, xtick={1,10,20,30,40,49}]
\addplot[clbase,dashed,thick,domain=0:50,samples=2]{69.2};
\addplot[clmodel,thick,mark=*,mark size=0.8pt] coordinates {
(1,18.9)(2,53.8)(3,55.8)(4,61.5)(5,56.4)(6,63.9)(7,62.7)(8,60.3)(9,65.0)(10,61.8)
(11,60.8)(12,59.6)(13,59.8)(14,61.9)(15,58.7)(16,59.8)(17,60.6)(18,61.1)(19,60.2)(20,55.8)
(21,53.3)(22,53.5)(23,53.0)(24,50.7)(25,50.2)(26,50.5)(27,50.7)(28,47.9)(29,49.0)(30,47.8)
(31,46.6)(32,47.0)(33,48.0)(34,48.1)(35,47.2)(36,46.7)(37,48.3)(38,49.0)(39,47.6)(40,46.9)
(41,47.4)(42,46.9)(43,46.9)(44,47.9)(45,47.4)(46,47.4)(47,47.9)(48,48.1)(49,48.1)};
\node[star,star points=5,fill=clgain,draw=none,inner sep=1.6pt] at (axis cs:9,65.0) {};
\node[font=\scriptsize,anchor=south west,text=clgain] at (axis cs:9,66) {peak, epoch 9};
\node[font=\scriptsize,anchor=south east,text=clbase] at (axis cs:49,70) {baseline 69.2\,\%};
%
% ---------- (b) D2, filtered, 14 epochs ----------
\nextgroupplot[title={\footnotesize (b) \sfilt, 14 epochs},
               xmin=0, xmax=15, xtick={1,4,...,14}]
\addplot[clbase,dashed,thick,domain=0:15,samples=2]{79.0};
\addplot[clmodel,thick,mark=*,mark size=1pt] coordinates {
(1,56.1)(2,72.2)(3,73.2)(4,73.5)(5,71.8)(6,73.8)(7,71.8)(8,71.8)(9,69.8)(10,70.8)
(11,69.1)(12,67.6)(13,68.2)(14,67.2)};
\node[star,star points=5,fill=clgain,draw=none,inner sep=1.6pt] at (axis cs:6,73.8) {};
\node[font=\scriptsize,anchor=south,text=clgain] at (axis cs:6,75) {epoch 6};
\node[font=\scriptsize,anchor=south east,text=clbase] at (axis cs:14.5,80) {79.0\,\%};
%
% ---------- (c) D3, mobile users, 9 epochs ----------
\nextgroupplot[title={\footnotesize (c) \smob, 9 epochs},
               xmin=0, xmax=10, xtick={1,3,...,9}]
\addplot[clbase,dashed,thick,domain=0:10,samples=2]{43.5};
\addplot[clmodel,thick,mark=*,mark size=1pt] coordinates {
(1,19.8)(2,40.2)(3,42.3)(4,42.5)(5,43.0)(6,42.3)(7,39.9)(8,38.7)(9,38.4)};
\node[star,star points=5,fill=clgain,draw=none,inner sep=1.6pt] at (axis cs:5,43.0) {};
\node[font=\scriptsize,anchor=south,text=clgain] at (axis cs:5,44.5) {epoch 5};
\node[font=\scriptsize,anchor=north east,text=clbase] at (axis cs:9.7,41.5) {43.5\,\%};
%
% ---------- (d) D3 + hour tokens, 11 epochs ----------
\nextgroupplot[title={\footnotesize (d) \smob + hour tokens, 11 epochs},
               xmin=0, xmax=12, xtick={1,3,...,11}]
\addplot[clbase,dashed,thick,domain=0:12,samples=2]{45.5};
\addplot[clmodel,thick,mark=*,mark size=1pt] coordinates {
(1,21.1)(2,39.0)(3,41.9)(4,41.1)(5,42.2)(6,41.3)(7,40.0)(8,39.4)(9,38.1)(10,38.8)(11,37.1)};
\node[star,star points=5,fill=clgain,draw=none,inner sep=1.6pt] at (axis cs:5,42.2) {};
\node[font=\scriptsize,anchor=south,text=clgain] at (axis cs:5,43.6) {epoch 5};
\node[font=\scriptsize,anchor=north east,text=clbase] at (axis cs:11.7,43.5) {45.5\,\%};
\end{groupplot}
\end{tikzpicture}}}
\caption{Validation accuracy per epoch on four datasets}
\label{fig:validation-curves}
\figtext{%
The same shape, four times, on four different datasets. In each panel the
horizontal axis is the training epoch and the vertical axis is top-1 accuracy
on validation users the run never trains on; the solid line is the model, the
star its best epoch, the dashed line that dataset's persistence baseline.
Validation climbs for four to nine epochs, peaks, then decays for tens of
epochs while the training loss keeps falling. In panel~(a) the peak is
65.0\,\% at epoch~9 against 48.1\,\% at the last epoch, so reading the end of
a run instead of its best point costs sixteen points and forty epochs of GPU
time. Note that the dashed line moves between panels, 69.2\,\% on the
unfiltered day and 43.5\,\% once the corpus is restricted to mobile users,
because the users themselves are different; the panels are comparable only on
the distance between the two lines, and in all four that distance is negative.
Per-epoch values are recovered by digitising each run's own validation curve,
with the vertical axis calibrated on the two values printed on that curve, so
an individual point carries about $\pm1\pt$ of reading error; the figures
quoted in the text are the reported ones. Qwen2.5-0.5B + QLoRA throughout.}
\figsource{\nbllm}{train\_cellid\_llm.ipynb}{the training loop whose per-epoch validation scores these curves are}
\end{figure}


This principle corresponds to \textbf{\panelref{fig:three-rules}{a}}. Fine-tuning language models on small trajectory datasets leads to severe overfitting when trained for excessive epochs. As shown in \textbf{Figure~\ref{fig:validation-curves}}, validation accuracy on unseen held-out subscribers reaches a peak within 4 to 9 training passes before declining, even as training loss continues to decrease.

In an extended 49-epoch test run shown in \textbf{\panelref{fig:validation-curves}{a}}, training loss decreased continuously, whereas held-out validation accuracy peaked at epoch 9 (65.0\% validation accuracy, 64.2\% test accuracy) before degrading to 48.1\% (48.7\% test accuracy) at epoch 49. Evaluating at final epoch completion would incur a \textbf{15.5 percentage point loss} in test accuracy on held-out subscribers. Implementing early stopping at peak validation accuracy alongside regularization (dropout 0.1, weight decay, learning rate scheduling) preserved model generalization while reducing training time from 72 to 29 minutes. \textbf{Figures~\panelnum{fig:validation-curves}{b}}, \textbf{\panelnum{fig:validation-curves}{c}}, and \textbf{\panelnum{fig:validation-curves}{d}} confirm that this overfitting dynamic holds consistently across dataset splits. All subsequent reported model metrics evaluate the optimal checkpoint saved at peak validation accuracy.

\FloatBarrier

\subsection{Principle 2: Context Fraction and Evaluation Split Sensitivity}

% T3 -- The context/target split swept on two datasets (slides 27 and 37).
\begin{table}[!htbp]
\centering\small
\caption{Moving the cut moves the ruler with it}
\label{tab:context-sweep}
\begin{tabular}{@{}lrrrrrrc@{}}
\toprule
\textbf{past / rest} & \textbf{Valid.} & \textbf{Model} & \textbf{Rule} & \textbf{Distance} & \textbf{top-3} & \textbf{top-5} & \textbf{Beats it?} \\
\midrule
\multicolumn{8}{@{}l}{\textit{\sfilt}, three separate trainings, one per cut}\\
\addlinespace[2pt]
70 / 30          & 72.3\,\% & 76.5\,\% & 76.8\,\% & $-0.3$ & 87.1\,\% & 90.1\,\% & no \\
\textbf{80 / 20} & 73.8\,\% & \textbf{80.1\,\%} & 79.0\,\% & $\mathbf{+1.1}$ & 89.1\,\% & 91.6\,\% & \textbf{yes} \\
90 / 10          & 75.7\,\% & 77.7\,\% & 80.1\,\% & $-2.4$ & 87.2\,\% & 90.5\,\% & no \\
\midrule
\multicolumn{8}{@{}l}{\textit{\smob}, one training at 80/20, evaluated at seven cuts}\\
\addlinespace[2pt]
30 / 70 & 43.0\,\% & 41.5\,\% & 42.9\,\% & $-0.9$ & 62.1\,\% & 69.9\,\% & no \\
40 / 60 & 43.0\,\% & 42.0\,\% & 43.2\,\% & $-0.6$ & 62.1\,\% & 69.7\,\% & no \\
50 / 50 & 43.0\,\% & 42.0\,\% & 43.3\,\% & $-0.7$ & 62.0\,\% & 69.4\,\% & no \\
60 / 40 & 43.0\,\% & 42.7\,\% & 44.0\,\% & $-0.9$ & 62.1\,\% & 69.5\,\% & no \\
70 / 30 & 43.0\,\% & 43.6\,\% & 44.3\,\% & $\mathbf{-0.3}$ & 62.1\,\% & 69.5\,\% & no \\
80 / 20 & 43.0\,\% & 44.6\,\% & 45.5\,\% & $-0.9$ & 63.9\,\% & 70.7\,\% & no \\
90 / 10 & 43.0\,\% & 47.6\,\% & 50.9\,\% & $-2.1$ & 68.0\,\% & 74.3\,\% & no \\
\bottomrule
\end{tabular}

\figtext{%
\figpara{Table~\ref{tab:context-sweep}: parameters and column descriptions.}{Only where each person's day
is cut into the part the model reads and the part it must produce: ``70 / 30''
means the model sees the first 70\,\% and is graded on the last 30\,\% in Table~\ref{tab:context-sweep}.
\textbf{Valid.} is the score on people held back to decide when to stop training,
\textbf{Model} the score on the test people, \textbf{Rule} what ``repeat the last
antenna'' scores on those very same moments, and \textbf{Distance} the model minus
the rule, which is the only column that means anything.}

\figpara{Table~\ref{tab:context-sweep}: persistence baseline variation across context fractions.}{The \textbf{Rule} column in Table~\ref{tab:context-sweep} changes from row
to row, because moving the cut changes which antenna counts as ``the last one
seen''. On \sfilt{} exactly one cut out of three clears the rule, by 1.1, while
the other two lose ground by 0.3 and 2.4; one favourable cell out of three is not
a result, and the next iteration made it disappear. On \smob{} no cut clears the
rule at all: the model climbs steadily from 41.5 to 47.6\,\% as it is given more
history, which looks like progress until one notices the rule climbing just as
fast. Each \sfilt{} row is a separate training run; the \smob{} block is one model
read seven times, and the giveaway is its identical \textbf{Valid.} column.}}

\figsource{\nbllm}{train\_cellid\_llm.ipynb}{the context-fraction sweeps}
\end{table}


This principle corresponds to \textbf{\panelref{fig:three-rules}{b}}. The prefix context fraction (share of trajectory provided to the model prior to prediction) directly affects task difficulty and baseline performance. \textbf{Table~\ref{tab:context-sweep}} presents model performance across varying context fractions (30\% to 90\%) on \Dfilt{} and \Dmob{}. 

As shown in \textbf{Table~\ref{tab:context-sweep}}, moving the context fraction cut alters the persistence baseline score because it shifts which event counts as ``the most recent cell.'' On \Dfilt{}, evaluating at an 80\% context fraction yields 80.1\% model accuracy against a 79.0\% persistence baseline (+1.1\%), whereas evaluating at a 70\% context fraction trails the baseline by 0.3\%. On \Dmob{}, absolute model accuracy increases from 41.5\% to 47.6\% as the context fraction expands from 30\% to 90\%; however, baseline accuracy increases at an identical rate (42.9\% to 50.9\%), leaving the performance delta negative across all cuts. Consequently, apparent accuracy gains resulting strictly from expanding historical context represent measurement ruler artifacts rather than model improvements, requiring context fractions to be held constant during comparative evaluations.

\FloatBarrier

\subsection{Principle 3: Baseline Upper Bounds and Task Room Constraints}

This principle corresponds to \textbf{\panelref{fig:three-rules}{c}}. On unfiltered daily logs, the persistence baseline achieves 69.2\% accuracy due to high stationary periods. A model achieving 64.4\% accuracy on such raw data has not failed in learning mobility patterns; rather, it has been evaluated on a benchmark dominated by stationary trivial predictions. Restricting benchmarks to active mobile trajectories (reducing baseline score to 45.5\%) creates a viable dynamic evaluation range.

\subsection{Methodological Implications for Hyperparameter Optimization}

Applying these three evaluation principles summarized in \textbf{Figure~\ref{fig:three-rules}} revealed that hyperparameter tuning alone yields limited performance improvements on uniform cohorts. Individual predictions revealed stark disparities: stationary trajectories were predicted with near 100\% confidence, whereas active transition trajectories resulted in incorrect baseline-like predictions. This variance across individual subscriber profiles motivated the individual user analysis presented in Chapter~\ref{ch:user}.

\FloatBarrier

\subsection{Principle 3: Baseline Upper Bounds and Task Room Constraints}

This principle is the one drawn in \textbf{\panelref{fig:three-rules}{c}}. On unfiltered daily logs, the persistence baseline achieves 69.2\% accuracy due to high stationary periods. A model achieving 64.4\% accuracy on such raw data has not failed in learning mobility patterns; rather, it has been evaluated on a benchmark dominated by stationary trivial predictions. Restricting benchmarks to active mobile trajectories (reducing baseline score to 45.5\%) creates a viable dynamic evaluation range.

\subsection{Methodological Implications for Hyperparameter Optimization}

Applying these three evaluation principles summarized in \textbf{Figure~\ref{fig:three-rules}} revealed that hyperparameter tuning alone yields limited performance improvements on uniform cohorts. Individual predictions revealed stark disparities: stationary trajectories were predicted with near 100\% confidence, whereas active transition trajectories resulted in incorrect baseline-like predictions. This variance across individual subscriber profiles motivated the individual user analysis presented in Chapter~\ref{ch:user}.

\section{Conclusion}
\label{sec:retro-tools}

This chapter established a standardized model fine-tuning architecture (Qwen2.5-0.5B fine-tuned via QLoRA with explicit cell and hour tokenization) and a rigorous evaluation protocol. Model capacity and PEFT mechanics remained fixed, allowing subsequent experimental gains to be attributed strictly to data representation and user conditioning mechanisms.

Establishing early stopping (\textbf{Figure~\ref{fig:validation-curves}}), standardized context fraction splits (\textbf{Table~\ref{tab:context-sweep}}), and active mobility cohort benchmarking eliminated false positive findings and provided a reliable foundation for subsequent research chapters.


